% - The two files: protein-level score matrix (1,853,074 proteins × 40 functional score columns, values in [0,1], NaN when 
%              protein not annotated with that function) and metadata file (same 1,853,074 rows, columns: accession, host, host\_type, 
%              split)\\
% - Scale of the dataset: 18,474 unique phages after aggregation\\
% - The split column: integer 0–4 encoding a genome-cluster-based partition — phages in the same genome cluster are always in the 
%   same fold, ensuring test phages are genuinely different from training phages at the sequence level\\
% - Class distribution: gram-neg vs gram-pos vs unknown\\ 
% - The NaN structure:  NaN is structural missingness, not noise — a protein that was not annotated with a given function simply has 
%   no score for that function, which is different from a score of zero



\chapter{Data and Preprocessing Pipeline}
\label{sec4: data}
\noindent In this section, we describe the dataset architecture \ref{sec4.1: data description}, the mathematical aggregation 
operators used to combine protein-level functional annotations into phage-level score vectors \ref{sec4.2: aggregation}, and the 
masking strategy developed to eliminate bias caused by cross-fold protein cluster contamination \ref{sec4.3: contamination fix}.

\section{Data Description and Structural Missingness}
\label{sec4.1: data description}
Our dataset comprises of $N_{\text{prot}} = 1\,853\,074$ protein sequences extracted from $18\,474$ unique bacteriophages. 
The base models use the EMPATHI pipeline (Section~\ref{sec2.2.2: plms}) to predict functional annotations across categories such 
as lysin, tail sheath, portal protein, and capsid \cite{coclet2021global}. An entry $s_{p,m} \in [0,1]$ is the model's predicted 
probability that protein $p$ carries function $m$. In this work we treat these scores as fixed features in our conformal framework 
without retraining the underlying classifiers.

Our pipeline operates on four primary data structures:
\begin{enumerate}[label=(\roman*)]
    \item \textbf{Protein Score Matrix:} Matrix of raw predictions $\mathbf{S} \in ([0,1] \cup \{\text{NA}\})^{N_{\text
          {prot}} \times M}$, where entry $s_{p,m}$ denotes the predicted probability that protein $p$ carries function $m$. For 
           the Gram-type task, we evaluate the $K$ functional categories aggregated across all hosts; for the host-level task, 
           $N \times K = 2\,160$ host-function pairs.
    \item \textbf{Phage Metadata:} Mappings that link each protein ID to its accession, host genus label $H \in \mathcal{H}$ 
          ($N = 54$ genera), host Gram-type $G \in \{\text{gram-neg}, \text{gram-pos}\}$, and protein cluster identifier (\texttt{pc}).
    \item \textbf{Cross-Validation Partition:} Integer split assignments $f \in \{0, 1, 2, 3, 4\}$ assigned at the genome level, 
           they ensure that all proteins from a given phage share the same split.
    \item \textbf{Protein Split Mask Matrix:} A binary assignment matrix mapping each protein to its evaluation split across all 
           host-function models.
\end{enumerate}

The dataset contains some structural missingness: approximately $30.6\%$ to $36.5\%$ of protein category scores are 
missing ($\text{NA}$) because phages possess only a subset of functional categories. An entry $s_{p,m} = \text{NA}$ indicates that 
protein $p$ was not annotated with function $m$. This is deterministic rather than stochastic: a capsid protein cannot carry a 
tail-fiber score, and an unannotated cluster does not provide any predictions. 

Formally, we can say that there exists a missingness mask $\mathbf{M} \in \{0,1\}^{N_{\text{prot}} \times 
M}$ such that:
\begin{equation}
\label{eq: structural-missingness}
s_{p,m} = \text{NA} \iff M_{p,m} = 0.
\end{equation}

This missingness reflects the biological identity of the phage, so it is not missing randomly. Imputing $\text{NA}$ 
values to zero would falsely treat unannotated proteins as evidence against infection which would violate the nonconformity 
interpretations of Section~\ref{sec3.3: nc transformation}. We therefore preserve missingness by using NaN aware aggregation 
(Section~\ref{sec4.2: aggregation}).

\begin{table}[htbp]
\centering
\begin{tabular}{lccc}
\toprule
\textbf{Set} & \textbf{Total Phages} & \textbf{Gram-Negative} & \textbf{Gram-Positive} \\
\midrule
Full Dataset & 18,474 & 10,735 & 7,708 \\
Test Split (Fold 0) & 4,433 & 2,811 & 1,611 \\
Validation Split (Fold 1) & 3,173 & -- & -- \\
Training Splits (Folds 2--4) & 10,868 & 7,924 & 6,097 \\
\bottomrule
\end{tabular}
\caption{Dataset Summary and Class Distributions}
\label{tab: dataset summary}
\end{table}

